% =====================================================================
% Paper 10: The middle-space (Morrow, Silver \& De Paoli)
% PERSONA_CATEGORY: Classic HCD design persona (empirically grounded design artifact)
% PERSONA_SUBCATEGORY: Persona-as-LLM-co-created design deliverable (human--LLM workflow from real interview data)
% PERSONA_SUBCATEGORY: Persona-as-requirements-bridge (input to scenarios and functionality prioritisation)
% =====================================================================

\subsection*{Paper 10 --- \emph{The middle-space: an evolution of AI-collaborative design practices for HCD design element generation through the lens of research discovery} (Morrow, Silver \& De Paoli)}

This paper is the clearest instance so far of the \emph{classic} Human-Centred Design understanding of persona: personas are ``HCD tools that help designers and developers create user-centred systems by representing lived experiences and user expectations --- including emotions, thoughts, feelings, views, and actions,'' delivered together with accompanying user scenarios and grounded in Cooper's and Nielsen's canonical formulations \cite{cooper1999, nielsen1990}. The persona is therefore neither a simulated respondent nor a conversational agent, but a \emph{narrative design artefact} representing a group of real users, produced as a project deliverable and used for downstream design decision-making.

What is novel is not the concept but its production pipeline. The authors position their contribution in a ``middle space'' where qualitative research practice coexists with LLM-driven processing, and explicitly reject synthetic data as a replacement for real user data: personas must remain anchored in authentic, contextually rich human participation. Operationally, personas are the third of five sequential steps: (1) 57 semi-structured interviews with researchers-as-users and industry stakeholders across two EU projects; (2) LLM-assisted thematic analysis (GPT-4o-mini on Azure via the authors' open-source \textsc{TALLMesh} tool), prompted with Braun and Clarke steps 1--4 and constrained to surface themes aligned with Nielsen's components of a basic persona; (3) LLM generation of ten personas and matching scenarios from researcher-selected themes against a predefined structural template; (4) further prompting of the personas/scenarios into candidate system functionalities, traced back to user data via Sankey diagrams; and (5) a proposed (not yet executed) longitudinal reproducibility test of artefact stability as LLMs evolve.

Validity is treated as a socially negotiated, human-in-the-loop property rather than a statistical one. Drafts were reviewed in a half-day interpretive team workshop, manually corrected against the empirical data, then evaluated anonymously by project consortium partners using a survey adapted from the Technology Acceptance Model. Reported quality failures are internal-coherence errors of the persona-as-document kind (a persona described as a 22-year-old with 20 years' research experience, corrected to 42; domain-inappropriate tool mentions), not failures of behavioural fidelity. Persona here is thus operationalized as a \emph{reproducible, auditable deliverable} whose warrant comes from traceability to interview data, team reflexivity and stakeholder endorsement --- an epistemically pragmatic account concerned with trust, transparency and ``how-to'' methodological guidance rather than with agentic simulation.

\begin{table*}[t]
\centering
\caption{Running taxonomy of persona conceptualizations across workshop papers (updated through Paper 10).}
\label{tab:persona-taxonomy}
\small
\begin{tabular}{p{0.5cm} p{3.6cm} p{3.4cm} p{3.6cm} p{3.6cm}}
\toprule
\textbf{\#} & \textbf{Paper} & \textbf{Main category} & \textbf{Subcategory} & \textbf{Operationalization / unit} \\
\midrule
1 & AI Personas for UX Research in Large Scale Retail (Sharma et al., Walmart) &
Surrogate user persona (LLM-simulated participant) &
(a) Persona-as-test-participant; (b) Persona-as-latent-trait-vector (activation steering) &
Demographic/behavioural archetype profiles prompted into a multimodal LLM to produce synthetic usability feedback; validated against 150 human participants. Secondary: persona as steering direction in BLIP-2 activations \\
\addlinespace
2 & AI Personas in Highly Regulated Environments: Tools for Calibrating User Reliance? (Vasiliu \& Guarese) &
System-facing agent persona (deployed AI character) &
Persona-as-interaction-metaphor for trust/reliance calibration (incl.\ adaptive / metaphor-fluid personas) &
Two prompt-encoded conversational personas (\textit{Expert Operator} vs.\ \textit{Machine}) defining tone, stance and dialogue structure of a voice CAI; manipulated as an experimental variable and evaluated via trust/workload/usability measures and a behavioural disclosure probe \\
\addlinespace
3 & Evaluating Conversational Agent Integration in Peer Support Groups (Uetova et al.) &
Surrogate user persona (LLM-simulated participant) &
Persona-as-data-derived individual profile (behavioural clone) for multi-agent social simulation; explicitly contrasted with large ``persona catalogues'' (NLP persona datasets) &
One persona per real past participant: LLM-generated summary of that person's chat history, demographics and survey answers (tone, emoji use, thematic focus), used to condition synthetic peer replies in a replayed group chat; combined with probabilistic responder selection and ``buddy'' pairing; framed as a pre-deployment testbed for agent prompt/logic refinement, with stated limits on fidelity and inclusivity \\
\addlinespace
4 & Grounding Persona-Driven Usability Simulation in Established Standards (Touir, Barika Ktata \& Soui) &
Surrogate user persona (LLM-simulated participant) &
Persona-as-autonomous-task-agent (embodied evaluator in a perceive--decide--act loop); Persona-as-edge-case/accessibility probe &
Supervisor-Agent auto-synthesises personas from UI purpose + source code using a constrained schema (goal, tech-literacy level, accessibility constraint); each persona conditions a VLM/LLM User-Agent that clicks, scrolls and types on a live prototype, emitting action traces, error logs and think-aloud self-reports mapped to ISO 9241-110, Nielsen's heuristics and WCAG. Named exemplars (``Martha R.,'' ``David L.'') act as edge-case probes; governance layer forbids unsupported demographic stereotyping and requires provenance and human escalation \\
\addlinespace
5 & Ideal Persona AI in Real-time (Pham Thi Ngoc Anh) &
Self-expression persona (user's idealized self rendered by GenAI) &
(a) Persona-as-generated-self-image (real-time text-to-image self-portrait); (b) Persona-as-ethical-behaviour-envelope (transparent ``default ethical persona'' of the system) &
Persona defined as ``a personality belonging to a participant'' --- an ideal self co-authored with GenAI. Six designers wrote ideal-persona descriptions with ChatGPT, then fed five keywords/questions as prompts to StreamDiffusion in TouchDesigner and watched a live image of that persona; screen capture, field notes, think-aloud and post-study questions analysed thematically. Breakdowns: inconsistent persona prototypes, uncanny valley, disconnection/trust erosion. Proposes moving persona design from ``creative writing'' to ``structural engineering'': prompt training, declared default ethical persona, ethical constraints exposed as UI/UX controls, continuous user control in real time \\
\addlinespace
6 & Mapping Personas to Text Transformations: A Taxonomy Outline for Content Adaptation (Sillano, De Russis, Cal\`o, Troncy \& Lisena) &
Content-adaptation audience persona (NLP prompt-conditioning target reader) &
(a) Persona-as-decomposable-trait-vector mapped to auditable text-transformation operators; (b) Persona-as-target-reader-profile for content personalization &
Persona = specification of the intended reader that conditions an LLM rewrite, critiqued as a ``black-box prompt modifier''. Decomposed into four textually-correlated dimensions (demographic, domain knowledge, cognitive, contextual) synthesised from role-playing/personalization surveys; each dimension mapped in a taxonomy table to named operators (lexical, structural, content, style, tone) with worked examples. Two tabular persona profiles (P1 novice student, P2 science journalist) drive a five-step operator chain rewriting a thermodynamics sentence with GPT-5.4, each edit traceable to an operator + persona dimension. Goal: transparency, bias auditability and writer agency; mappings not yet empirically validated \\
\addlinespace
7 & \textsc{Synonymix}: Privacy-Preserving Life Story Personas via Graph-Based Abstraction (Chen \& Kumar) &
Surrogate user persona (LLM-simulated participant) &
(a) Persona-as-narrative-life-story profile (high-fidelity behavioural proxy for a real individual); (b) Persona-as-privacy-preserving synthetic composite (graph-sampled from many real people); (c) Persona-as-collective artifact (``unigraph'' for meso-level sensemaking) &
Persona framed as data grounding for a generative agent, on a fidelity continuum from ``flattened'' demographic personas to McAdams-style life story personas. Each persona is converted into a typed knowledge graph (Subject / Factual / Interpretive nodes, with I-nodes produced by LLM ``expert reflection''), genericised and semantically deduplicated into a merged \emph{unigraph}; new personas are sampled by a Thematic Random Walk anchored on an interpretive node. Evaluated as three 30-agent banks (Demographic, Life Story, Synonymix) answering 108 GSS 2022 items, scored by EMD/TVD distributional distance ($p<0.001$, $r=0.59$ on ordinal items) and by Maximum Source Contribution ($\overline{MSC}=0.129$) as a structural privacy metric; DP over node histograms deferred. Persona additionally reframed as a shareable collective representation supporting cross-temporal sensemaking and consensus building rather than prediction \\
\addlinespace
8 & Temporal Stability of LLM Personas: A Prerequisite for Social Science Simulation (Gonnermann-Müller) &
Surrogate user persona (LLM-simulated participant) --- here broadened to \emph{simulated human research subject} &
(a) Persona-as-simulated-research-subject (social-science simulation agent); (b) Persona-as-psychometrically-measurable trait configuration (stability/drift construct); (c) Implicit split between persona-as-claimed-identity (self-report) and persona-as-enacted-behaviour (observer rating) &
Persona = prompt-encoded set of characteristics, behavioural tendencies and goals configuring an LLM agent to stand in for a human subject. Operationalized as four graded conditions on a clinical dimension (high / moderate / low ADHD + no-persona default), instantiated via three semantically equivalent prompt variants across seven LLMs. Stability measured with a multi-informant design: first-person self-report on the 12-item Conners' ADHD Index vs.\ blind ratings by three LLM observers scoring behaviour only. Experiment I: 50 independent runs per condition (between-conversation stability, small $SD$s, monotonic scaling with intensity). Experiment II: 18-turn dialogues probed at turns 6/12/18 --- self-reports stable, observer-rated expression declines $\sim$20\% with rising variance, attributed to context dilution and RLHF-driven normativity. Persona treated as a validity risk object: drift toward an ``average persona'' would make multi-agent emergence an artifact; proposes reporting standards (mid-session consistency checks, variability metrics, seed/config replication, full documentation) \\
\addlinespace
9 & The Bias Paradox: How AI Personas Can Overcome Human Limitations in UX Research (Celik) &
Surrogate user persona (LLM-simulated participant) --- framed as \emph{bias-resistant baseline} complementing human research &
(a) Persona-as-interviewable-research-informant (on-demand stand-in for recruited participants); (b) Persona-as-reactivated-design-artifact (classic static persona document turned conversational agent); (c) Persona-as-comparative-bias-profile (distinct, more predictable bias set than human participants) &
Six research-based personas distilled from $\sim$70 customer interviews (wealth level, investment style, goals, challenges, behavioural patterns) and instantiated as ChatGPT CustomGPT conversational agents made company-wide available so teams can consult user perspectives without scheduling research; the same personas also function conventionally as artifacts ``presented'' in a design-thinking workshop. Case study contrast: a 40-participant workshop with portfolio managers present in a luxury hotel yielded feedback judged contaminated by social desirability bias, context effects, relational anchoring and acquiescence; each bias is mapped to a corresponding persona ``immunity'' (no environmental perception, no relationships, no hospitality), while persona liabilities (training-data/configuration bias, demographic skew, sycophancy, hallucination) are argued to be consistent and detectable via adversarial prompting. Proposes a preliminary decision heuristic (when personas vs.\ humans are preferable) and raises the validation paradox of using possibly biased human responses as ground truth; notes organizational ``AI credibility'' concerns limiting adoption \\
\addlinespace
10 & The Middle-Space: An Evolution of AI-Collaborative Design Practices for HCD Design Element Generation (Morrow, Silver \& De Paoli) &
Classic HCD design persona (empirically grounded design artifact) --- explicitly Cooper/Nielsen lineage, contrasted with purely synthetic personas &
(a) Persona-as-LLM-co-created design deliverable (human--LLM ``middle space'' workflow over real qualitative data); (b) Persona-as-requirements-bridge feeding scenarios and functionality prioritisation; (c) Persona-as-reproducibility/validity object (traceability, longitudinal stability testing) &
Persona defined as an HCD tool representing lived experiences and expectations (emotions, thoughts, feelings, views, actions) plus paired user scenarios. Five-step pipeline: 57 semi-structured interviews across two EU research-discovery projects $\rightarrow$ LLM-assisted thematic analysis (GPT-4o-mini/Azure via open-source \textsc{TALLMesh}, Braun \& Clarke steps 1--4, themes aligned to Nielsen's basic-persona components) $\rightarrow$ LLM drafting of 10 personas + scenarios from a predefined template $\rightarrow$ prompt-derived functionality ideas traced to user data via Sankey diagrams $\rightarrow$ proposed longitudinal ``in-the-wild'' reproducibility testing. Human-in-the-loop validation via half-day interpretive team workshop, manual correction against transcripts, and anonymous consortium-partner survey adapted from the Technology Acceptance Model; reported errors are document-coherence flaws (22-year-old with 20 years' experience $\rightarrow$ 42; domain-mismatched tools). Position: synthetic data may support but must not replace authentic human participation; accountability in the ``middle space'' as basis for trust \\
\bottomrule
\end{tabular}
\end{table*}
